\documentclass[11pt,a4paper,sans]{moderncv} % Font sizes: 10, 11, or 12; paper sizes: a4paper, letterpaper, a5paper, legalpaper, executivepaper or landscape; font families: sans or roman

\moderncvstyle{classic} % CV theme - options include: 'casual' (default), 'classic', 'oldstyle' and 'banking'
\moderncvcolor{green} % CV color - options include: 'blue' (default), 'orange', 'green', 'red', 'purple', 'grey' and 'black'

\usepackage{lipsum} % Used for inserting dummy 'Lorem ipsum' text into the template
\usepackage{graphicx}
\usepackage[scale=0.75]{geometry} % Reduce document margins
\usepackage[official]{eurosym}
%\setlength{\hintscolumnwidth}{3cm} % Uncomment to change the width of the dates column
%\setlength{\makecvtitlenamewidth}{10cm} % For the 'classic' style, uncomment to adjust the width of the space allocated to your name

%----------------------------------------------------------------------------------------
%	NAME AND CONTACT INFORMATION SECTION
%----------------------------------------------------------------------------------------

\firstname{Johnny} % Your first name
\familyname{van Doorn} % Your last name

% All information in this block is optional, comment out any lines you don't need
%\title{Curriculum Vitae}
\address{Amsterdam}
%\mobile{0683695554}
\email{j.b.vandoorn@uva.nl}
%\website{https://scholar.google.com/citations?user=7TL5ZKsAAAAJ&hl=en}


\photo[80pt][0.5pt]{photo2_lo} % The first bracket is the picture height, the second is the thickness of the frame around the picture (0pt for no frame)


%----------------------------------------------------------------------------------------

\begin{document}
\makecvtitle % Print the CV title

%----------------------------------------------------------------------------------------
%	EDUCATION SECTION
%----------------------------------------------------------------------------


\section{Teaching}

\subsection{Courses}
\cventry{2019-current}{Research Methods \& Statistics (12 EC)}{\textsc{University of Amsterdam}}{}{}{Introductory statistics course for first year Psychology Bachelor's students.}

\cventry{2023-current}{Scientific and Statistical  Reasoning (9 EC)}{\textsc{University of Amsterdam}}{}{}{Intermediate statistics course for second year Psychology Bachelor's students.}

\cventry{2024-current}{Bayesian Statistics (6 EC)}{\textsc{University of Amsterdam}}{}{}{Specialization course for second year Psychological Methods Bachelor's students.}

\cventry{2021-2023}{Basic Skills in Mathematics, Programming \& Statistics (6 EC)}{\textsc{University of Amsterdam}}{}{}{Specialization course for second year Psychological Methods Bachelor's students.}

\cventry{2021-2023}{Statistiek voor de Levenswetenschappen(6 EC)}{\textsc{University of Amsterdam}}{}{}{Introductory course in statistics for second year B\`{e}ta-Gamma Bachelor's students.}

\cventry{2017-2020}{Programming in Psychological Science (6 EC)}{\textsc{University of Amsterdam}}{}{}{High intensity programming course in R, part of the Psychology Research Master's program.}


\subsection{Books}
\cvitem{2025}{Field, A., van Doorn, J., \& Wagenmakers, E.-J., \textit{Discovering statistics using JASP}. London: Sage. Companion website: \href{https://discoverjasp.com}{discoverjasp.com}.}
\cvitem{2023}{van Doorn, J., \textit{A Brief Introduction to Bayesian Inference}. Available at \href{https://johnnydoorn.github.io/IntroductionBayesianInference/}{johnnydoorn.github.io/IntroductionBayesianInference}.}

\subsection{Workshops}

\cventry{}{Annual Workshop -- Discovering Statistics with JASP}{\textsc{Various}}{}{}{Together with Eric-Jan Wagenmakers. See \href{https://jasp-stats.org/workshopdiscoveringstatistics/}{jasp-stats.org/workshopdiscoveringstatistics}.}

\cventry{}{Bayesian Inference in JASP}{\textsc{Various}}{}{}{Introductory workshops in Bayesian hypothesis testing and estimation using the statistics software JASP:}
\cvitem{2026}{OPAM workshop on Bayesian inference with JASP}
\cvitem{2026}{Workshop at St George's University}
\cvitem{2026}{UvA Library -- Machine Learning Basics}
\cvitem{2025}{Workshop at KU Leuven}
\cvitem{April 2023}{Full-day workshop for the TeaP society, held online}
\cvitem{March 2023}{Full-day workshop for the TeaP Conference held in Trier, Germany}
\cvitem{January 2023}{Full-day workshop for the TeaP society, held online}
\cvitem{August 2022}{Full-day workshop for the European Conference on Visual Perception held in Nijmegen, the Netherlands}
\cvitem{May 2022}{Half-day workshop for Hogrefe Publishers held in Amsterdam, the Netherlands}
\cvitem{June 2019}{Half-day workshop for Interlearn/Oefenweb held in Amsterdam, the Netherlands}
\cvitem{June 2019}{Webinar for the Society for Personality and Social Psychology - with Alexander Etz and Julia Haaf}
\cvitem{April 2019}{Three day Stan workshop for Leuven University, held in Leuven, Belgium-  with Don van den Bergh and Fabian Dablander}
\cvitem{Feb 2019}{Half-day workshop for the Society for Personality and Social Psychology Conference, held in Portland, USA - with Alexander Etz and Quentin Gronau}
\cvitem{July 2018}{Half-day workshop for the Society for the Improvement of Psychological Science, held in Grand Rapids, USA}
\cvitem{March 2018}{Two day workshop for the Adam Mickiewicz University in, held in Poznan, Poland - with Quentin Gronau}
\cvitem{Dec 2017}{Two day workshop for the Postgraduate School for Experimental Psychopathology, held in Heeze, the Netherlands - with Quentin Gronau and Noah van Dongen}



\subsection{Organization}
\cventry{}{Annual Workshop - JAGS and WinBUGS Workshop:  Bayesian Modeling for Cognitive Science}{Amsterdam}{}{}{Part  of  the  organizing  committee}
\cvitem{2019}{August 26-30,  in Amsterdam, the Netherlands}
\cvitem{2018}{August 20-24,  in Amsterdam, the Netherlands }
\cvitem{2017}{August 21-25,  in Amsterdam, the Netherlands }
\cvitem{2016}{August 15-19,  in Amsterdam, the Netherlands }

\cventry{}{Annual Workshop - Theory and Practice of Bayesian Hypothesis Testing}{Amsterdam}{}{}{Part  of  the  organizing  committee}
\cvitem{2019}{August 22-23,  in Amsterdam, the Netherlands }
\cvitem{2018}{August 27-28,  in Amsterdam, the Netherlands }
\cvitem{2017}{August 28-29,  in Amsterdam, the Netherlands}
\cvitem{2016}{August 22-23,  in Amsterdam, the Netherlands}

~\\
\cventry{2019}{PMLabs}{Amsterdam}{}{}{Lab meeting for the whole department of Psychological Methods, where all staff members show each other, and the students, what they are working on.}
~\\

\subsection{Supervision}
\cventry{}{Bachelor thesis}{}{}{}{}
\cventry{}{Master thesis}{}{}{}{}
\cventry{}{BDS internship}{}{}{}{}
\cventry{}{BDS thesis}{}{}{}{}
\cventry{}{Honours internship}{}{}{}{}

\newpage

~\\


\subsection{Miscellaneous}
\cventry{2016-current}{JASP Team Member}{}{\url{https://jasp-stats.org/}}{}{JASP is an open-source statistics software that aims to increase the accessibility, transparency, and intuitiveness of conducting statistical analyses. My work has focused on adding features to the ANOVA, regression, and $t$-test. Recently, my function has focused more on supervising the software transition from SPSS to JASP, developing teaching materials, and helping out other teachers using JASP.}

\cventry{2020-current}{Study Advisor}{}{University of Amsterdam}{}{Active study advisor for the Psychological Methods group, for Bachelor (psychology) and Master (BDS) students. Monitoring students' well-being, organizing information sessions, advising on electives/study direction, and coordinating the student delegacy.}

\cventry{2022-current}{JASP Community Coordinator}{}{JASP/University of Amsterdam}{}{Together with E--J Wagenmakers, setting up an inter-university collaboration around the open-source statistics software JASP. The  \href{https://jasp-stats.org/community-vision-and-goals}{goal of the foundation}  is to ensure the financial longevity of the software, and to promote open software in statistics education.}

\cventry{2022-current}{Bachelor Curriculum Coordinator Psychological Methods}{}{University of Amsterdam}{}{Part of the committee that revises the educational targets of the Psychology bachelor's program (in Dutch: ``Project Zichtbare Leerlijnen").}

\cventry{2022-current}{Coordinating Transition to JASP}{}{University of Amsterdam}{}{The bachelor curriculum in psychology is transition from SPSS to JASP in its statistics courses. My tasks consist of making sure this transition goes smoothly, making sure every department is supported and offering workshops/tutorial sessions.}

\cventry{2019-2020}{PhD Representative}{Department of Psychological Methods}{University of Amsterdam}{}{Duties include organizing lab meetings, offering counseling to junior researchers, and being a spokesperson at policy meetings.}

\cventry{2019-2020}{Member of the Bachelor's Program Reform Committee}{Department of Psychological Methods}{University of Amsterdam}{}{Together with eight colleagues, we are working on a reform of the Psychological Methods Bachelor program, in order to accommodate structural changes imposed by the Faculty of Psychology.}


\end{document}